\section*{Capítulo \ref{ch:b1:continuity} --- Limites e Continuidade}

\begin{solution}{exo:b1:continuity:1}
Tome $u_n = \frac{1}{2\pi n + \pi/2}$ e $v_n = \frac{1}{2\pi n}$:
ambas tendem a $0^+$ e, no entanto, $\sin\frac{1}{u_n} = 1$ e
$\sin\frac{1}{v_n} = 0$. Duas sequências, dois limites diferentes das
imagens: pelo \cref{thm:b1:continuity:seqchar}, não há limite em $0^+$.

$x \sin\frac1x$: comprimida por $\abs{x\sin\frac1x} \leq \abs x \to 0$,
de modo que o limite em $0$ existe e vale $0$.
\end{solution}

\begin{solution}{exo:b1:continuity:2}
$f = \lfloor\cdot\rfloor$ é \omterm{def:b1:continuity:continuous}{contínua} em $\R \setminus \Z$
(localmente constante) e descontínua em cada $n \in \Z$: limite à
esquerda $n - 1$, valor $n$.

$g(x) = x - \lfloor x\rfloor$: mesmos pontos de descontinuidade (a
identidade é \omterm{def:b1:continuity:continuous}{contínua}, de modo que $g$ herda os saltos de $f$); em $n \in
\Z$, o limite à esquerda é $1 \neq 0 = g(n)$.

$h(x) = \lfloor x\rfloor + (x - \lfloor x\rfloor)^2$: em $\intco{n}{n
+ 1}$, $h(x) = n + (x - n)^2$, \omterm{def:b1:continuity:continuous}{contínua} ali; em $x = n$, o limite à
esquerda é $(n-1) + 1 = n = h(n)$: os saltos se cancelam. $h$ é \omterm{def:b1:continuity:continuous}{contínua}
em $\R$ (e estritamente crescente).
\end{solution}

\begin{solution}{exo:b1:continuity:3}
$P(x) = x^5 - 3x + 1$: $P(-2) = -32 + 6 + 1 = -25 < 0$; $P(0) = 1 >
0$; $P(1) = -1 < 0$; $P(2) = 27 > 0$. Três mudanças de sinal nos
segmentos disjuntos $\intcc{-2}{0}$, $\intcc{0}{1}$, $\intcc{1}{2}$:
pelo \cref{thm:b1:continuity:ivt}, ao menos três raízes. (Sendo de
grau $5$, $P$ tem no máximo cinco; um estudo de variação mostraria
exatamente três.)
\end{solution}

\begin{solution}{exo:b1:continuity:4}
Seja $P = a_{2m+1}X^{2m+1} + \dots$ com $a_{2m+1} > 0$ (caso contrário, troque
$P$ por $-P$). Fatorando o termo dominante, $P(x) = a_{2m+1}
x^{2m+1}\bigl(1 + o(1)\bigr)$ quando $x \to \pm\infty$: assim, $P(x) \to
+\infty$ em $+\infty$ e $-\infty$ em $-\infty$. Tome $a$ com
$P(a) < 0$ e $b$ com $P(b) > 0$: o teorema do valor intermediário em
$\intcc{a}{b}$ fornece uma raiz.
\end{solution}

\begin{solution}{exo:b1:continuity:5}
Seja $g(x) = f(x) - x$, \omterm{def:b1:continuity:continuous}{contínua} em $\intcc{0}{1}$. Como $f$ leva
em $\intcc{0}{1}$: $g(0) = f(0) \geq 0$ e $g(1) = f(1) - 1 \leq
0$. Pelo \cref{thm:b1:continuity:ivt}, $g(c) = 0$ para algum $c$: um
ponto fixo.

Necessidade das hipóteses: em $\intcc{0}{1}$, a \omterm{def:b1:logic:map}{aplicação} descontínua
$f(x) = 1$ para $x \leq \frac12$, $f(x) = 0$ para $x > \frac12$, não tem
ponto fixo; no \omterm{prop:b1:reals:intervals}{intervalo} $\intoo{0}{1}$ (que não é segmento), $f(x) =
\frac x2$ é \omterm{def:b1:continuity:continuous}{contínua} em $\intoo{0}{1}$, sem ponto fixo
(o candidato $0$ está faltando); em $\R$, $f(x) = x + 1$.
\end{solution}

\begin{solution}{exo:b1:continuity:6}
$g(x) = \cos x - x$ é \omterm{def:b1:continuity:continuous}{contínua}, $g(0) = 1 > 0$, $g(1) = \cos 1 -
1 < 0$: existe uma solução em $\intoo{0}{1}$
(\cref{thm:b1:continuity:ivt}). Unicidade: $g$ é estritamente
decrescente em $\R$ --- para $x \leq 0$, $g(x) \geq 1 - x > 0$ não tem
zero, de qualquer modo; e $g'(x) = -\sin x - 1 \leq 0$, com igualdade apenas em
pontos isolados ($x \equiv -\frac\pi2 \bmod 2\pi$), de sorte que $g$ é
estritamente decrescente (\cref{ch:b1:derivative}; alternativamente: em
$\intcc{0}{1}$, $\cos$ é estritamente decrescente e $-x$ também, logo $g$
é). Uma função estritamente monótona se anula no máximo uma vez.
\end{solution}

\begin{solution}{exo:b1:continuity:7}
Fixe $M = f(0) + 1$. Existe $A > 0$ com $f(x) \geq M$ para $\abs x
\geq A$ (definição dos dois limites infinitos; tome o maior
limiar). No segmento $\intcc{-A}{A}$, o teorema de Weierstrass
(\cref{thm:b1:continuity:evt}) fornece $c$ com $f(c) =
\inf_{\intcc{-A}{A}} f \leq f(0)$. Para $\abs x \geq A$: $f(x) \geq M
> f(0) \geq f(c)$. Logo, $f(c)$ é o mínimo global.
\end{solution}

\begin{solution}{exo:b1:continuity:8}
No segmento $\intcc{0}{T}$, $f$ é limitada e atinge as suas cotas
(\cref{thm:b1:continuity:evt}); pela periodicidade, essas são as cotas
em todo o $\R$, ainda atingidas.

Seja $g(x) = f(x + \frac T2) - f(x)$, \omterm{def:b1:continuity:continuous}{contínua}. Então
\[
g(0) + g\bigl(\tfrac T2\bigr)
= \bigl(f(\tfrac T2) - f(0)\bigr) + \bigl(f(T) - f(\tfrac T2)\bigr)
= f(T) - f(0) = 0 :
\]
$g(0)$ e $g(\frac T2)$ têm sinais opostos (ou um deles se anula), de modo
que o teorema do valor intermediário em $\intcc{0}{T/2}$ dá $c$ com
$g(c) = 0$, isto é, $f(c + \frac T2) = f(c)$.
\end{solution}

\begin{solution}{exo:b1:continuity:9}
Primeiro a desigualdade: para $0 \leq y \leq x$,
\[
\bigl(\sqrt y + \sqrt{x - y}\bigr)^2 = x + 2\sqrt{y(x-y)} \geq x,
\]
de modo que $\sqrt x \leq \sqrt y + \sqrt{x - y}$, isto é, $\sqrt x - \sqrt y
\leq \sqrt{x - y}$. Portanto, $\abs{\sqrt x - \sqrt y} \leq
\sqrt{\abs{x - y}}$ para todos $x, y \geq 0$.

\omterm{def:b1:continuity:uniform}{Continuidade uniforme}: dado $\varepsilon > 0$, tome $\delta =
\varepsilon^2$; então $\abs{x - y} \leq \delta$ implica $\abs{\sqrt x
- \sqrt y} \leq \sqrt\delta = \varepsilon$.

Não lipschitziana perto de $0$: $\frac{\sqrt x - \sqrt 0}{x - 0} =
\frac{1}{\sqrt x} \to +\infty$ quando $x \to 0^+$, de sorte que constante $k$ alguma
pode dominar todos os quocientes de diferenças.
\end{solution}

\begin{solution}{exo:b1:continuity:10}
Suponha $f$ \omterm{def:b1:logic:inj}{injetiva}, \omterm{def:b1:continuity:continuous}{contínua} e não estritamente monótona.
Então existem $a < b < c$ em $I$ com $f(b)$ não entre $f(a)$ e
$f(c)$ --- com efeito, se, para \emph{todas} as triplas, o valor do
meio estivesse entre os das pontas, $f$ seria monótona (compare dois
pares quaisquer; uma curta verificação por casos). Digamos que $f(b) > \max(f(a), f(c))$ (o outro caso é simétrico: troque $f$ por $-f$). Escolha $v$ com
$\max(f(a),
f(c)) < v < f(b)$. Pelo teorema do valor intermediário aplicado em
$\intcc{a}{b}$ e em $\intcc{b}{c}$, existem $u_1 \in
\intoo{a}{b}$ e $u_2 \in \intoo{b}{c}$ com $f(u_1) = v = f(u_2)$:
dois pontos distintos com imagens iguais, contradizendo a injetividade.
\end{solution}

\begin{solution}{exo:b1:continuity:11}
$f(0) = f(0+0) = 2f(0)$ dá $f(0) = 0$; $f(-x) = -f(x)$ vem de $0 =
f(x - x)$. Ponha $\alpha = f(1)$. Indução: $f(n) = n\alpha$ para $n
\in \N$, e depois para $n \in \Z$, pela imparidade. Para $q \in \N^*$:
$q\,f(\frac pq) = f(p) = p\alpha$ (some $\frac pq$ a si mesmo $q$
vezes), de modo que $f(\frac pq) = \alpha\frac pq$: $f = \alpha\,
\mathrm{id}$ em $\Q$.

Sejam agora $x \in \R$ e $(r_n)$ uma sequência de racionais com $r_n \to
x$ (\omterm{def:b1:topology:dense}{densidade}, \cref{thm:b1:reals:density}, aplicada em
\omterm{prop:b1:reals:intervals}{intervalos} encaixados; ou $r_n = \frac{\lfloor nx\rfloor}{n}$). \omterm{def:b1:continuity:continuous}{Continuidade}:
$f(x) = \lim f(r_n) = \lim \alpha r_n = \alpha x$.
\end{solution}

\begin{solution}{exo:b1:continuity:12}
Seja $\varepsilon > 0$. Pelo limite em $+\infty$, existe $A$ com
$\abs{f(x) - \ell} \leq \frac\varepsilon2$ para $x \geq A$; logo, para
$x, y \geq A$: $\abs{f(x) - f(y)} \leq \varepsilon$ (sem necessidade
de proximidade).

No segmento $\intcc{0}{A + 1}$, o teorema de Heine
(\cref{thm:b1:continuity:heine}) dá $\delta_0 > 0$ para esse
$\varepsilon$; ponha $\delta = \min(\delta_0, 1)$.

Tome agora quaisquer $x, y \geq 0$ com $\abs{x - y} \leq \delta$, digamos $x
\leq y$. Se $y \leq A + 1$: os dois estão no segmento, e o
$\delta_0$ de Heine se aplica. Caso contrário, $y > A + 1$ e, então, $x \geq y - 1 >
A$: os dois estão em $\intco{A}{+\infty}$, onde o argumento do limite
se aplica. Nos dois casos, $\abs{f(x) - f(y)} \leq \varepsilon$: \omterm{def:b1:continuity:uniform}{continuidade
uniforme}.
\end{solution}

\begin{solution}{pb:b1:continuity:1}
\textbf{1.} Para $n \in \N$: $f(nx) = n f(x)$, por indução
($f((n+1)x) = f(nx) + f(x)$). Além disso, $f(0) = 2f(0)$ dá $f(0) =
0$, e $0 = f(x - x) = f(x) + f(-x)$ dá a imparidade, de modo que $f(nx) =
nf(x)$ para $n \in \Z$. Para $r = \frac pq$: $q\,f\bigl(\tfrac
pq x\bigr) = f(px) = p\,f(x)$, logo $f(rx) = r f(x)$: $f$ é
$\Q$-linear.

\textbf{2.} A aditividade dá, para todos $x$ e $h$:
\[
f(x + h) - f(x) = f(h) = f(x_0 + h) - f(x_0) .
\]
Quando $h \to 0$, o lado direito tende a $0$ pela \omterm{def:b1:continuity:continuous}{continuidade} em
$x_0$; portanto, $f(x + h) \to f(x)$: \omterm{def:b1:continuity:continuous}{continuidade} em todo $x$. Então o
\cref{exo:b1:continuity:11} fornece $f(x) = f(1)\,x$.

\textbf{3.} Duas funções aditivas \omterm{def:b1:continuity:continuous}{contínuas} são da forma
$cx$ e $c'x$ (questão 2); se elas coincidem no \omterm{def:b1:structures:subgroup}{subgrupo} \omterm{def:b1:topology:dense}{denso}
$\Z + \sqrt2\,\Z$ (\cref{exo:b1:reals:9}), então $c\,g = c'g$ para
algum $g \neq 0$ nele: $c = c'$, e as funções são iguais. Sem
\omterm{def:b1:continuity:continuous}{continuidade}: a $\Q$-linearidade só amarra valores em
$\Q$-combinações, e $1, \sqrt2$ são $\Q$-independentes ($\sqrt2
\notin \Q$), de modo que $f(1) = 0$, $f(\sqrt2) = 1$ é compatível e
força, no \omterm{def:b1:structures:subgroup}{subgrupo},
\[
f(m + n\sqrt2) = m\,f(1) + n\,f(\sqrt2) = n .
\]

\textbf{4.} Para $t \in \intcc{0}{\ell}$: $a + t \in
\intcc{a}{b}$, de modo que $f(t) = f(a + t) - f(a) \leq M - f(a) =: M'$.

\textbf{5.} Para $t \in \intcc{0}{\ell}$, também $\ell - t \in
\intcc{0}{\ell}$, e a aditividade dá $f(t) = f(\ell) - f(\ell
- t) \geq f(\ell) - M'$. Portanto, $\abs{f} \leq C$ em
$\intcc{0}{\ell}$, com $C = \max\bigl(\abs{M'}, \abs{f(\ell) -
M'}\bigr)$.

\textbf{6.} Para $t \in \intcc{0}{\ell/n}$: $nt \in
\intcc{0}{\ell}$ e $f(t) = \frac{f(nt)}{n}$ (questão 1), de modo que
$\abs{f(t)} \leq \frac Cn$. Dado $\varepsilon > 0$, escolha $n >
\frac C\varepsilon$: para $0 \leq h \leq \frac\ell n$,
$\abs{f(h)} \leq \varepsilon$, e, para $h$ negativo, use $f(h) =
-f(-h)$. Assim, $f(h) \to 0 = f(0)$ quando $h \to 0$: \omterm{def:b1:continuity:continuous}{continuidade} em $0$,
logo em toda parte (questão 2), logo $f(x) = f(1)x$.

\textbf{7.} Se $f$ é não decrescente em $\intcc{a}{b}$, então
$f(a) \leq f(x) \leq f(b)$ ali: limitada, logo linear pela questão
6, $f(x) = cx$; e $c(b - a) = f(b) - f(a) \geq 0$ força $c
\geq 0$.

\textbf{8.} (a)$\Rightarrow$(b)$\Rightarrow$(c): trivial.
(c)$\Rightarrow$(a): questão 2. (a)$\Rightarrow$(d): uma função linear
é monótona em toda parte. (d)$\Rightarrow$(e): uma função monótona
em $\intcc{a}{b}$ é limitada ali pelos seus valores nas
extremidades. (e)$\Rightarrow$(a): questões 4--6. Os cinco
enunciados são equivalentes --- a escada da regularidade.

\textbf{9.} A questão 4 usou apenas a \omterm{def:b1:reals:bounds}{cota superior} $M$; a questão 5
\emph{deduziu} a cota inferior da superior via a
reflexão $f(t) = f(\ell) - f(\ell - t)$; a questão 6 rodou então sobre
$\abs f \leq C$. Assim: aditiva e limitada \emph{superiormente} num
\omterm{prop:b1:reals:intervals}{intervalo} não degenerado já implica linear. Para uma cota
inferior, aplique isso a $-f$ (aditiva, limitada superiormente). Degrau
(e) mais forte: uma cota unilateral num \omterm{prop:b1:reals:intervals}{intervalo} basta.

\textbf{10.} Se $\frac{f(x)}{x}$ fosse uma mesma constante $c$ para todo
$x \neq 0$, $f$ seria linear; assim, existem $u, v$ não nulos
com $\frac{f(u)}u \neq \frac{f(v)}v$, isto é, $u f(v) - v f(u)
\neq 0$: o determinante dos vetores $(u, f(u))$, $(v,
f(v))$ é não nulo, e eles geram $\R^2$.

\textbf{11.} Para $r, s \in \Q$: $f(ru + sv) = r f(u) + s f(v)$
(questão 1 duas vezes mais aditividade), de modo que o gráfico contém
\[
\bigl(ru + sv,\; r f(u) + s f(v)\bigr)
= r\,(u, f(u)) + s\,(v, f(v)) .
\]
Dados $(x_0, y_0)$ e $\varepsilon > 0$: o sistema real $2 \times 2$
$a(u, f(u)) + b(v, f(v)) = (x_0, y_0)$ tem uma (única)
solução $(a, b)$, pois o determinante é não nulo. Tome
racionais $r_n \to a$, $s_n \to b$: então $r_n(u, f(u)) +
s_n(v, f(v)) \to (x_0, y_0)$ coordenada a coordenada, e cada um desses
pontos está no gráfico: o gráfico é \omterm{def:b1:topology:dense}{denso} em $\R^2$.

\textbf{12.} Sejam $I$ um \omterm{prop:b1:reals:intervals}{intervalo} não degenerado, $x_0$ o seu
ponto médio e $M$ arbitrário: a \omterm{def:b1:topology:dense}{densidade} fornece um ponto do gráfico a
menos de $\min\bigl(\frac{\abs I}{2}, 1\bigr)$ de $(x_0, M + 1)$, isto é,
$x \in I$ com $f(x) > M$: ilimitada em $I$, logo (questão 8)
descontínua em todo ponto e monótona em \omterm{prop:b1:reals:intervals}{intervalo} nenhum; e,
para qualquer alvo $y_0$, pontos do gráfico perto de $(x_0, y_0)$ dão $f(x)$
arbitrariamente próximo de $y_0$ com $x \in I$: $f(I)$ é \omterm{def:b1:topology:dense}{denso} em
$\R$. Isto é a questão 8 lida ao contrário: como todos os degraus são
equivalentes, uma função aditiva não linear deve falhar em \emph{todos}
eles, em toda parte.

\textbf{13.} Bem definida: $m + n\sqrt2 = m' + n'\sqrt2$ força
$(n - n')\sqrt2 = m' - m \in \Z$, de modo que $n = n'$ (caso contrário, $\sqrt2 \in
\Q$) e $m = m'$. A aditividade em $G$ é então clara, coordenada a
coordenada. Ilimitação perto de $0^+$: fixe $\varepsilon \in
\intoo{0}{1}$ e $N \in \N$. Para cada $n$ fixado com $\abs n
\leq N$, a condição $m + n\sqrt2 \in \intoo{0}{1}$ prende $m$
dentro de um \omterm{prop:b1:reals:intervals}{intervalo} de comprimento $1$: no máximo um inteiro $m$ por
$n$, de modo que no máximo $2N + 1$ elementos de $G \cap \intoo{0}{1}$ têm
$\abs{f} \leq N$. Mas $G \cap \intoo{0}{\varepsilon}$ é
infinito ($G$ é \omterm{def:b1:topology:dense}{denso}, \cref{exo:b1:reals:9}); logo, ele
contém algum $g$ com $\abs{f(g)} > N$: $f$ é ilimitada em
toda \omterm{def:b1:topology:open}{vizinhança} à direita de $0$. Estender $f$ a uma função aditiva
em $\R$ significa escolher valores coerentemente numa família de
reais que seja $\Q$-linearmente independente e gere $\R$ sobre $\Q$
--- uma base de Hamel; produzir uma exige o axioma da escolha,
e fórmula explícita alguma o consegue: construtivamente, só possuímos
o monstro em $G$.

\textbf{14.} $f(x) = f(\frac x2 + \frac x2) = f(\frac x2)^2 \geq
0$. Se $f(x_0) = 0$, então $f(x) = f(x - x_0)f(x_0) = 0$ para todo
$x$: excluído. Logo, $f > 0$ e $g = \ln \circ f$ é \omterm{def:b1:continuity:continuous}{contínua}
(\cref{prop:b1:continuity:algebra}), com $g(x + y) = g(x) +
g(y)$: pelo \cref{exo:b1:continuity:11}, $g(x) = cx$, de modo que $f(x) =
\eu^{cx}$. Reciprocamente, cada $\eu^{cx}$ serve: os
\omterm{def:b1:structures:morphism}{morfismos} \omterm{def:b1:continuity:continuous}{contínuos} $(\R, +) \to (\R^*, \times)$ são exatamente as
exponenciais.

\textbf{15.} $g(u) = f(\eu^u)$ é \omterm{def:b1:continuity:continuous}{contínua} e $g(u + v) =
f(\eu^u \eu^v) = g(u) + g(v)$: $g(u) = cu$, e todo $x > 0$
se escreve $x = \eu^u$ com $u = \ln x$: $f(x) = c\ln x$.

\textbf{16.} $h = \ln \circ f$ é \omterm{def:b1:continuity:continuous}{contínua} em
$\intoo{0}{+\infty}$ com $h(xy) = h(x) + h(y)$: pela questão 15,
$h(x) = c\ln x$, de modo que $f(x) = \eu^{c\ln x} = x^c$.

\textbf{17.} $x = y = 0$: $f(0) = 2f(0) + f(0)^2$, de modo que $f(0)(1 +
f(0)) = 0$. Se $f(0) = -1$: fazendo $y = 0$, $f(x) = f(x) +
f(0) + f(x)f(0) = -1$ para todo $x$: a constante $f \equiv -1$
(que de fato satisfaz a equação). Caso contrário, $f(0) = 0$; $h
= 1 + f$ é \omterm{def:b1:continuity:continuous}{contínua}, $h(0) = 1$, e
\[
h(x + y) = 1 + f(x) + f(y) + f(x)f(y) = h(x)\,h(y) :
\]
pela questão 14, $h(x) = \eu^{cx}$, isto é, $f(x) = \eu^{cx} - 1$
(o caso $c = 0$ dando $f \equiv 0$). Lista completa: $f
\equiv -1$ e $f(x) = \eu^{cx} - 1$, com $c \in \R$.

\textbf{18.} $x = y = 0$: $2f(0) = 4f(0)$, de modo que $f(0) = 0$. $x =
0$: $f(y) + f(-y) = 2f(y)$, de sorte que $f$ é par. $y = x$: $f(2x) =
4f(x)$.
Indução usando $(x, y) \to (nx, x)$:
\[
f((n{+}1)x) = 2f(nx) + 2f(x) - f((n{-}1)x)
= (2n^2 + 2 - (n-1)^2) f(x) = (n+1)^2 f(x) .
\]
Então $f(x) = f\bigl(q\cdot\frac xq\bigr) = q^2 f\bigl(\frac
xq\bigr)$ dá $f\bigl(\frac pq x\bigr) = \frac{p^2}{q^2}f(x)$:
$f(r) = r^2 f(1)$ em $\Q$ (a paridade trata os sinais). As duas
funções \omterm{def:b1:continuity:continuous}{contínuas} $f$ e $x \mapsto f(1)x^2$ coincidem no
\omterm{def:b1:logic:sets}{conjunto} \omterm{def:b1:topology:dense}{denso} $\Q$, logo em toda parte (questão 24): $f(x) =
f(1)\,x^2$; e todo $cx^2$ satisfaz a equação.

\textbf{19.} $g = f - f(0)$ é \omterm{def:b1:continuity:continuous}{contínua}, $g(0) = 0$, e
satisfaz a equação de Jensen (as constantes se cancelam). Fazendo $y =
0$: $g\bigl(\frac x2\bigr) = \frac{g(x)}{2}$. Então, para todos $x,
y$:
\[
\frac{g(x + y)}{2} = g\Bigl(\frac{x+y}{2}\Bigr)
= \frac{g(x) + g(y)}{2} ,
\]
de modo que $g$ é aditiva e \omterm{def:b1:continuity:continuous}{contínua}: $g(x) = cx$ (questão 2), e
$f(x) = cx + d$ com $d = f(0)$. Todas as funções afins satisfazem
Jensen: a lista está completa.

\textbf{20.} Indução em $n$. Para $n = 0$: $\lambda \in \{0,
1\}$, trivial. Suponha a desigualdade para todos os pesos
$\frac{k}{2^n}$. Um peso $\lambda = \frac{k}{2^{n+1}}$ com $k$
par se reduz ao nível $n$; para $k = 2j + 1$, $\lambda$ é o
ponto médio de $\lambda_1 = \frac{j}{2^n}$ e $\lambda_2 =
\frac{j+1}{2^n}$. Com $z_i = \lambda_i x + (1 - \lambda_i)y$:
$\lambda x + (1-\lambda)y = \frac{z_1 + z_2}{2}$, de modo que
\[
f\bigl(\lambda x + (1{-}\lambda)y\bigr)
\leq \frac{f(z_1) + f(z_2)}{2}
\leq \frac{(\lambda_1 + \lambda_2)f(x) + (2 - \lambda_1 -
\lambda_2)f(y)}{2}
= \lambda f(x) + (1 - \lambda)f(y) .
\]

\textbf{21.} Fixe $x, y$. As aplicações $\lambda \mapsto f(\lambda x +
(1 - \lambda)y)$ e $\lambda \mapsto \lambda f(x) + (1 -
\lambda)f(y)$ são \omterm{def:b1:continuity:continuous}{contínuas} em $\intcc{0}{1}$ (composição e
álgebra, \cref{prop:b1:continuity:algebra}). A desigualdade
vale nos pesos diádicos, que são \omterm{def:b1:topology:dense}{densos} em $\intcc{0}{1}$
(\cref{exo:b1:reals:8}); para $\lambda$ arbitrário, tome diádicos
$\lambda_n \to \lambda$ e passe ao limite
(\cref{thm:b1:continuity:seqchar} e \cref{thm:b1:seq:order}):
a desigualdade de convexidade vale para todo $\lambda \in
\intcc{0}{1}$ --- convexidade no ponto médio mais \omterm{def:b1:continuity:continuous}{continuidade} é
igual a convexidade (a noção do \cref{ch:b1:derivative}).

\textbf{22.} Uma $f$ aditiva não linear satisfaz
$f\bigl(\frac{x+y}{2}\bigr) = \frac{f(x) + f(y)}{2}$ exatamente
(questão 1 com $r = \frac12$, e depois aditividade): ela é
convexa no ponto médio, e até afim no ponto médio. Se ela satisfizesse a
desigualdade de convexidade plena em algum \omterm{prop:b1:reals:intervals}{intervalo} $\intcc{x}{y}$, então,
para $\lambda \in \intcc{0}{1}$: $f(\lambda x + (1-\lambda)y) \leq
\max(f(x), f(y))$ --- limitada superiormente num \omterm{prop:b1:reals:intervals}{intervalo} não degenerado,
logo linear pela questão 9: contradição. Assim, a \omterm{def:b1:continuity:continuous}{continuidade} na
questão 21 não é luxo: sem ela, a convexidade no ponto médio
controla apenas o esqueleto diádico enumerável, e o \omterm{def:b1:continuity:continuous}{contínuo}
entre eles corre solto.

\textbf{23.} $p(x) = \frac{x^2}{2}$ satisfaz $p(x+y) = p(x) +
p(y) + xy$. Se $f$ é uma solução \omterm{def:b1:continuity:continuous}{contínua} qualquer, $g = f - p$ é
\omterm{def:b1:continuity:continuous}{contínua} e aditiva, de modo que $g(x) = cx$:
\[
f(x) = \frac{x^2}{2} + cx , \qquad c \in \R ,
\]
e cada uma delas é solução: a lista está completa.

\textbf{24.} Princípio geral: se $u, v$ são \omterm{def:b1:continuity:continuous}{contínuas} e
coincidem num \omterm{def:b1:logic:sets}{conjunto} \omterm{def:b1:topology:dense}{denso} $D \subseteq \R$, então, para $x \in \R$, escolha
$d_n \in D$ com $d_n \to x$
(\cref{prop:b1:topology:closureprops}); $u(x) = \lim u(d_n) =
\lim v(d_n) = v(x)$. Seja agora $f$ \omterm{def:b1:continuity:continuous}{contínua} e aditiva no
\omterm{def:b1:structures:subgroup}{subgrupo} \omterm{def:b1:topology:dense}{denso} $G$. Fixe $x, y \in \R$ e tome $g_n \to x$,
$g'_n \to y$ com $g_n, g'_n \in G$; então $g_n + g'_n \to x + y$
e, pela \omterm{def:b1:continuity:continuous}{continuidade} sequencial em $x + y$, em $x$ e em $y$:
\[
f(x + y) = \lim f(g_n + g'_n) = \lim\bigl(f(g_n) +
f(g'_n)\bigr) = f(x) + f(y) :
\]
$f$ é aditiva em todo o $\R$ (logo linear, pela questão 2).

\textbf{25.} (i) Para $f$ aditiva: linear $\iff$ \omterm{def:b1:continuity:continuous}{contínua}
$\iff$ \omterm{def:b1:continuity:continuous}{contínua} num ponto $\iff$ monótona em algum \omterm{prop:b1:reals:intervals}{intervalo}
$\iff$ limitada (mesmo de um só lado) em algum \omterm{prop:b1:reals:intervals}{intervalo}. (ii) A
\omterm{def:b1:topology:dense}{densidade} do gráfico no plano mostra que a falha não é um defeito
local, e sim uma explosão global: acima de todo subintervalo os
valores se espalham por todo o $\R$, de modo que toda propriedade de
regularidade falha em toda parte de uma vez. (iii) A pescaria: $cx$,
$\eu^{cx}$, $c\ln x$, $x^c$, $cx^2$ e $cx + d$ --- seis
caracterizações e um método: transporte a equação para a de Cauchy,
demonstre o esqueleto em $\Q$ por indução, e promova a $\R$ por
\omterm{def:b1:topology:dense}{densidade} mais \omterm{def:b1:continuity:continuous}{continuidade}. (iv) A \omterm{def:b1:topology:dense}{densidade} de $\Q$ (ou dos diádicos)
carregou as promoções no \cref{exo:b1:continuity:11} e na questão 21; ela
nada carregou na questão 13, porque, sem \omterm{def:b1:continuity:continuous}{continuidade}, os
valores não se propagam de um \omterm{def:b1:logic:sets}{conjunto} \omterm{def:b1:topology:dense}{denso} ao seu \omterm{def:b1:topology:closure}{fecho} ---
a \omterm{def:b1:topology:dense}{densidade} transfere informação apenas ao longo da \omterm{def:b1:continuity:continuous}{continuidade}.
\end{solution}
